% =====================================================================
%  Relativistic modification of Chapter II, §15
%  Exact solutions of the characteristic value problem
%  for relativistic thermal (Bénard) instability
% =====================================================================
%
%  Classical reference: Chandrasekhar, Ch. II §15, eqs. (187)–(227)
%  Relativistic framework: BDNK causal first-order dissipation
%  (Bemfica et al. 2018; Kovtun 2019)
%  Agent: 9 (subagent for Bénard, part 4 — exact solutions)
% =====================================================================

\section{Exact solutions of the relativistic characteristic value problem}
\label{sec:rel-15}

We return to the relativistic eigenvalue problem formulated in the
preceding sections. The classical analysis of \S\ref{sec:15} determined
the critical Rayleigh number $R_c$ for three canonical boundary
configurations. Here we solve the corresponding relativistic problem,
where the eigenvalue equation is modified by two physical effects:
\begin{enumerate}
\item[(i)] the \emph{pressure contribution to inertia}
  $(\edensity + p)/c^{2}$, which replaces the rest-mass density
  $\rdensity$ in the inertial terms; and
\item[(ii)] the \emph{BDNK causal first-order heat transport}
  (Bemfica et al.\ 2018; Kovtun 2019), which ensures finite-speed
  thermal signal propagation through frame-dependent transport
  coefficients rather than relaxation times.
\end{enumerate}

\subsection{The relativistic eigenvalue equation}
\label{sec:rel-15-eigen}

In the Boussinesq limit of the relativistic perturbation equations
(cf.\ the development in the preceding relativistic sections), the
sixth-order eigenvalue equation for the vertical velocity perturbation
$W(z)$ becomes
\begin{equation}
\label{eq:rel-15-eigen}
\bigl(D^{2}-a^{2}\bigr)^{3}W = -\Rarel\,a^{2}\,W,
\end{equation}
where the relativistic Rayleigh number is
\begin{equation}
\label{eq:rel-15-Radef}
\Rarel = \frac{\alpha\,g\,d^{3}\,\Delta T}{\nu\,\kappa_{T}}
  \cdot\frac{1}{1 + p_{0}/(\edensity_{0}\,c^{2})}.
\end{equation}
Here $\edensity_{0}$ and $p_{0}$ are the equilibrium energy density
and pressure, and $\kappa_{T}$ is the thermal diffusivity. We define
the dimensionless relativistic parameter
\begin{equation}
\label{eq:rel-15-params}
\Xi \;\equiv\; \frac{p_{0}}{\edensity_{0}\,c^{2}},
\end{equation}
so that
\begin{equation}
\label{eq:rel-15-Rarel}
\Rarel = \frac{R}{1+\Xi},
\end{equation}
where $R$ is the classical Rayleigh number. In the non-relativistic
limit $\Xi \to 0$, one recovers $\Rarel \to R$.

\emph{Note:} In the BDNK first-order theory (Bemfica et al.\ 2018;
Kovtun 2019), there is no thermal relaxation time $\tauq$, so the
Israel--Stewart parameter $\Upsilon = \tauq\kappa_T/d^2$ does not
appear.  At marginal stability ($\sigma = 0$), the BDNK and
Israel--Stewart theories yield \emph{identical} results, since all
time-derivative terms vanish and only the first-order constitutive
relation $\Pi^{\mu\nu} = -2\shearvisc\,\sigma^{\mu\nu}$ contributes.

\begin{relcorrection}
\textbf{Key physical insight.}
The factor $(1+\Xi)^{-1}$ arises because the relativistic inertia of
the fluid is $(\edensity + p)/c^{2}$ rather than $\rdensity$. Since
$\edensity + p > \rdensity c^{2}$ for any fluid with non-zero pressure,
the effective inertia is \emph{larger}, and therefore $\Rarel < R$ for
the same temperature gradient. Equivalently, the pressure contribution
to inertia \emph{lowers} the effective Rayleigh number, making the
system \textbf{more unstable} at a given adverse temperature gradient.

In the BDNK formalism, no relaxation time $\tauq$ appears; causality
is instead ensured by the BDNK frame coefficients and coupled
transport-coefficient inequalities.  At the marginal state
($\sigma = 0$), these frame corrections do not contribute, and the
BDNK result coincides exactly with the Israel--Stewart result.
\end{relcorrection}

Since equation~\eqref{eq:rel-15-eigen} has the \emph{same operator
structure} as the classical equation~(\ref{eq:2-198}), the
eigenfunctions are unchanged for each boundary type. The entire
relativistic modification enters through the eigenvalue relation: one
replaces $R$ by $\Rarel$ everywhere. This crucial simplification means
that the classical solutions of Pellew \& Southwell, Low, and Reid \&
Harris carry over directly, with the critical Rayleigh number shifted
by the factor $(1+\Xi)$.


% -----------------------------------------------------------------
\subsection{Two free boundaries (relativistic)}
\label{sec:rel-15a}
% -----------------------------------------------------------------

The boundary conditions are identical to the classical case~(\ref{eq:2-187}):
\begin{equation}
\label{eq:rel-15-bc-free}
W = D^{2}W = D^{4}W = 0 \qquad\text{for } z = 0 \text{ and } 1.
\end{equation}
The eigenfunctions remain $W = A\sin n\pi z$ ($n = 1,2,\ldots$), and
substitution into~\eqref{eq:rel-15-eigen} gives the relativistic
characteristic equation
\begin{equation}
\label{eq:rel-15-Rfree}
\Rarel = \frac{(n^{2}\pi^{2}+a^{2})^{3}}{a^{2}}.
\end{equation}
For $n = 1$ the minimum over $a^{2}$ is attained at $a^{2} = \pi^{2}/2$
(the same critical wavenumber as in the classical theory), yielding
\begin{equation}
\label{eq:rel-15-Rcfree}
\boxed{\;
\Rarel^{(c)} = \frac{27}{4}\,\pi^{4} = 657.511.
\;}
\end{equation}
Expressing this in terms of the \emph{classical} Rayleigh number $R$
via~\eqref{eq:rel-15-Rarel}, the onset condition becomes
\begin{equation}
\label{eq:rel-15-Rcfree-classical}
R_{c}^{(\mathrm{rel})} = \frac{27}{4}\,\pi^{4}\,(1+\Xi)
  = 657.511\;\bigl[1 + \Xi\bigr].
\end{equation}

\begin{relcorrection}
The classical critical Rayleigh number $R_{c} = 27\pi^{4}/4 = 657.511$
is \emph{exact} in the relativistic theory provided one uses $\Rarel$
as the control parameter. When the onset condition is phrased in terms
of the classical (Newtonian) Rayleigh number $R$, the critical value
receives a multiplicative relativistic correction:
\[
R_{c}^{(\mathrm{rel})} = R_{c}^{(\mathrm{class})}
  \times (1+\Xi)
  = R_{c}^{(\mathrm{class})}
    \left[1 + \frac{p_{0}}{\edensity_{0}c^{2}}
    + \order{\Xi^2}\right].
\]
Since the factor $(1+\Xi) > 1$, a \emph{larger} classical
Rayleigh number is needed to trigger instability, but this is precisely
because the same physical temperature gradient corresponds to a
\emph{smaller} relativistic Rayleigh number (due to increased inertia).
Hence, for a given temperature gradient, the system is \textbf{more
unstable} in the relativistic regime.
\end{relcorrection}


% -----------------------------------------------------------------
\subsection{Two rigid boundaries (relativistic)}
\label{sec:rel-15b}
% -----------------------------------------------------------------

With the fluid confined between $z = \pm\tfrac{1}{2}$ and boundary
conditions
\begin{equation}
\label{eq:rel-15-bc-rigid}
W = DW = (D^{2}-a^{2})^{2}W = 0 \qquad\text{for } z = \pm\tfrac{1}{2},
\end{equation}
the operator structure is unchanged, and the general solution has the
same form as in the classical case (cf.\ equations~(\ref{eq:2-200})--(\ref{eq:2-206})),
with the parameter $\tau$ now defined through
\begin{equation}
\label{eq:rel-15-tau-rigid}
\Rarel\,a^{2} = \tau^{3}\,a^{6}.
\end{equation}
The even and odd transcendental equations~(\ref{eq:2-218})
and~(\ref{eq:2-225}) are therefore formally identical, and the
numerical solution of Reid and Harris applies verbatim. The critical
state for the lowest (even) mode is
\begin{equation}
\label{eq:rel-15-Rcrigid}
\boxed{\;
\Rarel^{(c)} = 1707.762, \qquad a_{c} = 3.117.
\;}
\end{equation}
In terms of the classical Rayleigh number:
\begin{equation}
\label{eq:rel-15-Rcrigid-classical}
R_{c}^{(\mathrm{rel})} = 1707.762\,(1+\Xi).
\end{equation}

For terrestrial laboratory conditions (e.g.\ water at room temperature),
$\Xi \sim p_{0}/(\rdensity c^{2}) \approx 10^{-10}$ so $R_{c}^{(\mathrm{rel})}$ is indistinguishable
from $1707.762$. The corrections become significant only for
relativistic fluids: in a quark--gluon plasma with
$p \approx \edensity/3$, one has $\Xi \approx 1/4$, giving
$R_{c}^{(\mathrm{rel})} \approx 1707.76 \times 1.25 = 2134.7$.

The eigenfunctions $W(z)$ and $(a^{2}\Rarel)^{-1}F(z)$ are
\emph{identical} to the classical expressions~(\ref{eq:2-220}) with
the same values of $q_{0}$, $q_{1}$, $q_{2}$
given in~(\ref{eq:2-221}), since these depend only on $a$ and $\tau$,
which take the same numerical values at criticality.


% -----------------------------------------------------------------
\subsection{One rigid and one free boundary (relativistic)}
\label{sec:rel-15c}
% -----------------------------------------------------------------

As in the classical case~(\S\ref{sec:15}(c)), the solution for one
rigid and one free boundary can be deduced from the odd solution for
two rigid boundaries applied to a layer of half the depth. The argument
carries over without modification to the relativistic case, and we
obtain
\begin{equation}
\label{eq:rel-15-Rcmixed}
\boxed{\;
\Rarel^{(c)} = 1100.65, \qquad a_{c} = 2.682.
\;}
\end{equation}
In terms of the classical Rayleigh number:
\begin{equation}
\label{eq:rel-15-Rcmixed-classical}
R_{c}^{(\mathrm{rel})} = 1100.65\,(1+\Xi).
\end{equation}


% -----------------------------------------------------------------
\subsection{Summary: classical versus relativistic critical values}
\label{sec:rel-15d}
% -----------------------------------------------------------------

Table~\ref{tab:rel-15-summary} collects the results for all three
boundary configurations.

\begin{table}[htbp]
\centering
\caption{Critical parameters for the onset of relativistic thermal
  instability compared with the classical (Newtonian) values.
  Here $\Xi = p_{0}/(\edensity_{0}c^{2})$.  In the BDNK first-order
  theory no IS relaxation parameter appears; at $\sigma=0$ the
  BDNK and IS results are identical.}
\label{tab:rel-15-summary}
\medskip
\begin{tabular}{@{}l c c c c@{}}
\toprule
Boundary type
  & $R_{c}^{(\mathrm{class})}$
  & $\Rarel^{(c)}$
  & $R_{c}^{(\mathrm{rel})}$
  & $a_{c}$ \\
\midrule
Both free
  & $657.511$
  & $657.511$
  & $657.511\,(1{+}\Xi)$
  & $2.221$ \\[4pt]
Both rigid
  & $1707.762$
  & $1707.762$
  & $1707.762\,(1{+}\Xi)$
  & $3.117$ \\[4pt]
One rigid, one free
  & $1100.65$
  & $1100.65$
  & $1100.65\,(1{+}\Xi)$
  & $2.682$ \\
\bottomrule
\end{tabular}
\end{table}

\noindent\textbf{Remarks.}
\begin{enumerate}
\item The critical wavenumber $a_{c}$ is \emph{unchanged} by the
  relativistic corrections, since the minimization of $\Rarel(a)$ yields
  the same functional form as in the classical theory.

\item The critical relativistic Rayleigh number $\Rarel^{(c)}$ coincides
  numerically with the classical $R_{c}$ in every case. All relativistic
  physics is absorbed into the \emph{definition} of $\Rarel$.

\item When the onset condition is expressed in terms of the classical
  Rayleigh number $R$, one finds $R_{c}^{(\mathrm{rel})} >
  R_{c}^{(\mathrm{class})}$ by the factor $(1+\Xi)$.
  Physically, this means that relativistic effects raise the
  \emph{apparent} threshold as measured by $R$, but the system is in
  fact \emph{more unstable} for a given temperature gradient because
  the enhanced inertia reduces $\Rarel$ below $R$.

\item For astrophysical applications (neutron-star oceans, accretion
  disks, quark--gluon plasma), $\Xi$ can be of order unity, producing
  $\mathcal{O}(1)$ shifts in the critical conditions.
\end{enumerate}


% -----------------------------------------------------------------
\subsection{Causality verification of critical modes}
\label{sec:rel-15-causality}
% -----------------------------------------------------------------

\begin{causalitycheck}
We must verify that the critical eigenmodes found above propagate
causally, i.e.\ no signal speed exceeds $c$.

\paragraph{Phase and group velocities of the critical mode.}
At marginal stability the perturbation is time-independent
($\sigma = 0$), so the critical mode is \emph{stationary}---it has
zero phase velocity. There is no propagation, and causality is
trivially satisfied.

\paragraph{Slightly supercritical regime.}
For $\Rarel$ slightly above $\Rarel^{(c)}$, the growth rate $\sigma$
is real and positive (exchange of stabilities). The corresponding
disturbance grows exponentially in time but does not propagate
horizontally faster than a diffusive speed. Specifically, the
perturbation envelope satisfies
\[
v_{\mathrm{signal}} \sim \sqrt{\kappa_{T}\,\sigma}
  \ll c,
\]
since both $\kappa_{T}$ and $\sigma$ are non-relativistic in magnitude
for the thermal instability problem.

\paragraph{Role of BDNK frame coefficients.}
In the BDNK first-order theory, causality is ensured by the choice of
hydrodynamic frame and the coupled inequalities on transport and frame
coefficients (Bemfica et al.\ 2018; Kovtun 2019; Hoult \& Kovtun 2020).
The characteristic speeds of the BDNK system are rigorously bounded by
$c$ when these inequalities are satisfied, without the need for
relaxation times $\tauq$ or $\taupi$.  At the marginal state
($\sigma = 0$), all time-derivative terms vanish, so the BDNK and
Israel--Stewart theories yield \emph{identical} critical Rayleigh
numbers and eigenfunctions.

\paragraph{Conclusion.}
All critical and near-critical modes of the relativistic B\'enard
problem propagate causally. The BDNK framework guarantees
that no acausal behaviour is introduced by the dissipative terms.
\end{causalitycheck}

% -----------------------------------------------------------------
\subsection{Astrophysical applications: quantitative critical Rayleigh numbers}
\label{sec:rel-15-applications}
% -----------------------------------------------------------------

\paragraph{Neutron star crust ($\Xi\approx 0.015$).}
$R_{c}^{(\mathrm{rel})} = 1707.76\times 1.015 = 1733.4$ (1.5\% increase).

\paragraph{Neutron star core ($\Xi\approx 0.14$).}
The APR equation of state gives $\Xi\approx 0.14$:
\begin{center}
\begin{tabular}{lccc}
\toprule
Boundary & $R_{c}^{(\mathrm{class})}$ & $R_{c}^{(\mathrm{rel})}$ & Shift \\
\midrule
Both free & 657.5 & 749.6 & $+14\%$ \\
Both rigid & 1707.8 & 1946.8 & $+14\%$ \\
Mixed & 1100.7 & 1254.7 & $+14\%$ \\
\bottomrule
\end{tabular}
\end{center}

\paragraph{Quark--gluon plasma ($\Xi\approx 1/3$).}
$R_{c}^{(\mathrm{rel})} = 1707.76\times 4/3 = 2277.0$ (rigid; 33\% increase).

\begin{figure}[htbp]
\centering
\includegraphics[width=0.95\textwidth]{fig_critical_Ra_table.pdf}
\caption{(a)~Critical Rayleigh numbers for 4 regimes and 3 boundary types.
(b)~Marginal stability curves $R_{c}(a)$ showing upward shift with $\xi$.}
\label{fig:critical-Ra-table}
\end{figure}

\section*{References}

\begin{enumerate}
\item F.~S.~Bemfica et al., \textit{Phys.\ Rev.\ D} \textbf{98}, 104064 (2018).
\item P.~Kovtun, \textit{JHEP} \textbf{10}, 034 (2019).
\item S.~Chandrasekhar, \textit{Hydrodynamic and Hydromagnetic Stability}, OUP (1961).
\item P.~Douchin and P.~Haensel, \textit{Astron.\ Astrophys.}\ \textbf{380}, 151 (2001).
\item A.~Akmal, V.~R.~Pandharipande, and D.~G.~Ravenhall, \textit{Phys.\ Rev.\ C} \textbf{58}, 1804 (1998).
\item P.~Romatschke and U.~Romatschke, \textit{Relativistic Fluid Dynamics}, CUP (2019).
\item L.~Rezzolla and O.~Zanotti, \textit{Relativistic Hydrodynamics}, OUP (2013).
\end{enumerate}
